\section{Random vectors}

A random variable is a real-valued measurable
function defined on a probability space
$(\Omega, {\cal F}, \mathrm{P})$ (when we talk of $\R$
as a measurable space, it will always be taken to mean
with the $\sigma$-algebra of Borel sets). Similarly,
we now wish to talk about \emph{vector}-valued measurable
functions on a probability space, i.e., functions
taking values in $\R^d$. First, we need to identify a
good $\sigma$-algebra of subsets of $\R^d$. Risking
some confusion, we will still call it the Borel
$\sigma$-algebra and denote it by ${\cal B}$, or
sometimes by ${\cal B}(\R^d)$.

\begin{defn}
\label{defi-borelrd}
The Borel $\sigma$-algebra on $\R^d$ is defined in
one of the following equivalent ways:

\begin{enumerate}[label=(\roman*)]
\item It is the $\sigma$-algebra generated by boxes of the
form

\[
(a_1,b_1) \times (a_2,b_2) \times \ldots \times (a_d,b_d).
\]

\item It is the $\sigma$-algebra generated by the balls in
$\R^d$.

\item It is the $\sigma$-algebra generated by the open sets
in $\R^d$.

\item It is the minimal $\sigma$-algebra of subsets of
$\R^d$ such that the coordinate functions
$\pi_i:\R^d\to\R$ defined by

\[
\pi_i({\bf x}) = x_i, \qquad i=1,2,\ldots,d
\]

\noindent are all measurable (where measurability is respect to
the Borel $\sigma$-algebra on the target space $\R$).
\end{enumerate}

\end{defn}

\begin{xexercise}
Check that the definitions above are indeed all
equivalent.
\end{xexercise}

\begin{defn}
A $\textbf{random ($d$-dimensional) vector}$ (or
$\textbf{vector random variable}$)
$\mathbf{X}=(X_1,X_2,\ldots,X_d)$ on a probability
space $(\Omega, {\cal F}, \mathrm{P})$ is a function
$\mathbf{X}:\Omega\to\R^d$ that is measurable (as a
function between the measurable spaces
$(\Omega, {\cal F})$ and $(\R^d, {\cal B})$.
\end{defn}

\begin{lem}
$\mathbf{X}=(X_1,\ldots,X_d)$ is a random vector if
and only if $X_i$ is a random variable for each
$i=1,\ldots,d$.
\end{lem}

\begin{proof}
If $\mathbf{X}$ is a random vector then each of its
coordinates $X_i = \pi_i \circ \mathbf{X}$ is a
composition of two measurable functions and therefore
(check!) measurable. Conversely, if $X_1,\ldots,X_d$
are random variables then for any box
$E=(a_1,b_1)\times (a_2,b_2)\times\cdots\times
(a_d,b_d) \subset \R^d$
we have

\[
\mathbf{X}^{-1}(E) = \cap_{k=1}^d X_i^{-1}((a_i,b_i)) \in {\cal F}.
\]

\noindent Therefore by \hyperref[defi-borelrd]{Definition~\ref*{defi-borelrd}} and \texorpdfstring{\hyperref[ex-checkrv]{Exercise~\ref*{ex-checkrv}}}{Exercise}, $\mathbf{X}$ is a random vector.
\end{proof}

\begin{xexercise}

\begin{enumerate}[label=(\roman*)]
\item Prove that any continuous function $f:\R^m\to\R^n$ is
measurable (when each of the spaces is equipped with
the respective Borel $\sigma$-algebra).

\item Prove that the composition $g\circ f$ of measurable
functions
$f:(\Omega_1, {\cal F}_1)\to (\Omega_2,{\cal F}_2)$
and
$g:(\Omega_2, {\cal F}_2)\to (\Omega_3, {\cal F}_3)$
(where $(\Omega_i, {\cal F}_i)$ are measurable spaces
for $i=1,2,3$) is a measurable function.

\item Deduce that the sum $X_1+\ldots+X_d$ of random
variables is a random variable.
\end{enumerate}

\end{xexercise}

\begin{xexercise}
Prove that if $X_1, X_2, \ldots$ is a sequence of
random variables (all defined on the same probability
space, then the functions

\[
\inf_n X_n, \quad \sup_n X_n, \quad \limsup_n X_n, \quad \liminf_n X_n
\]

\noindent are all random variables. Note: Part of the question
is to generalize the notion of random variable to a
function taking values in
$\overline{\mathbb{R}} = \R \cup \{-\infty, +\infty
\}$,
or you can solve it first with the additional
assumption that all the $X_i$'s are uniformly bounded
by some constant $M$.
\end{xexercise}
